\subsection{实值函数的代数运算}

设 \(f\) 与 \(g\) 是定义在集合 \(\mathbf{X}\) 上的两个实值函数；如果和 \(f(x) + g(x)\) {[}相应地，积 \(f(x)g(x)\) {]}对一切 \(x \in \mathbf{X}\) 有定义，则我们用 \(f + g\) （相应地，\(fg\) ）表示实值函数

\[
x \rightarrow f (x) + g (x) \quad [ \text { resp. } x \rightarrow f (x) g (x) ].
\]

同样，若 \(1 / f(x)\) 对一切 \(x\in \mathbf{X}\) 有定义，则 \(1 / f\) 表示函数 \(x\to 1 / f(x)\) 。

因此只要 \(f\) 不取值 \(0\) ，这最后一个函数便有定义；当 \(f\) 在区间 \([0, +\infty]\) （相应地，\([-\infty, 0]\) ）中取值时，\(1/f(x)\) 在约定 \(1/0 = +\infty\) （相应地，\(1/0 = -\infty\) ）之下便处处有定义；在这种情形函数 \(1/f\) 有定义。

设 X 由滤子 \(\mathfrak{F}\) 滤子化，且 \(\lim_{\mathfrak{F}} f\) 与 \(\lim_{\mathfrak{F}} g\) 存在。若一方面函数 \(f + g\) （相应地，\(fg\) ，\(1/g\) ）有定义，另一方面表达式 \(\lim_{\mathfrak{F}} f + \lim_{\mathfrak{F}} g\) （相应地，\(\lim_{\mathfrak{F}} f\) . \(\lim_{\mathfrak{F}} g\) ，\(1/\lim_{\mathfrak{F}} f\) ）有意义，则 \(\lim_{\mathfrak{F}} (f + g)\) {[}相应地，\(\lim_{\mathfrak{F}} fg\) ，\(\lim_{\mathfrak{F}} (1/f)\) {]}存在且等于该表达式，这是由于函数 \(x + y\) （相应地，\(xy\) ，\(1/x\) ）在其有定义的点处连续。

命题 12. 设 \(f\) 与 \(g\) 是定义在集合 \(\mathbf{X}\) 上的两个实值函数，并设 \(\mathbf{A}\) 是 \(\mathbf{X}\) 的一个非空子集。

\textit{(i) 我们有}

\begin{enumerate}
\def\labelenumi{(\arabic{enumi})}
\setcounter{enumi}{16}
\tightlist
\item
\end{enumerate}

\[
\sup _ {x \in \mathbf {A}} (f (x) + g (x)) \leqslant \sup _ {x \in \mathbf {A}} f (x) + \sup _ {x \in \mathbf {A}} g (x),\tag{17}
\]

\[
\sup _ {x \in \mathbf {A}} f (x) + \inf _ {x \in \mathbf {A}} g (x) \leqslant \sup _ {x \in \mathbf {A}} (f (x) + g (x)),\tag{18}
\]

只要这些不等式两边都有定义。

\begin{enumerate}
\def\labelenumi{(\roman{enumi})}
\setcounter{enumi}{1}
\tightlist
\item
  若 \(f(x)\) 与 \(g(x)\) 都 \(\geqslant 0\) （对每个 \(x \in \mathbf{A}\) ），则
\end{enumerate}

\begin{enumerate}
\def\labelenumi{(\arabic{enumi})}
\setcounter{enumi}{18}
\tightlist
\item
\end{enumerate}

\[
\sup _ {x \in \mathbf {A}} (f (x) g (x)) \leqslant \sup _ {x \in \mathbf {A}} f (x) \sup _ {x \in \mathbf {A}} g (x),\tag{19}
\]

\[
\sup _ {x \in \mathbf {A}} f (x) \inf _ {x \in \mathbf {A}} g (x) \leqslant \sup _ {x \in \mathbf {A}} (f (x) g (x))\tag{20}
\]

只要这些不等式两边都有定义。

\begin{enumerate}
\def\labelenumi{(\roman{enumi})}
\setcounter{enumi}{2}
\tightlist
\item
  若 \(f(x) \geqslant 0\) （对一切 \(x \in \mathbf{A}\) ），则
\end{enumerate}

\[
\sup _ {x \in \mathbf {A}} (1 / f (x)) = \inf _ {x \in \mathbf {A}} f (x)\tag{21}
\]

（约定 \(1 / 0 = +\infty\) ）。

设 \(\mathbf{H}\) 是 \(\mathbf{A}\) 的任意一个有限子集。若 \(x_0\) 是 \(\mathbf{H}\) 中使 \(f + g\) 取其最大值的点之一，则我们有

\[
f (x _ {0}) + g (x _ {0}) \leqslant \sup _ {x \in \mathbf {H}} f (x) + \sup _ {x \in \mathbf {H}} g (x);
\]

另一方面，若 \(x_{1}\) 是 \(H\) 中使 \(f\) 取其最大值的点之一，则

\[
f (x _ {1}) + g (x _ {1}) \geqslant \sup _ {x \in \mathbf {H}} f (x) + \inf _ {x \in \mathbf {H}} g (x);
\]

因此

\[
\sup _ {x \in \mathbf {H}} f (x) + \inf _ {x \in \mathbf {H}} g (x) \leqslant \sup _ {x \in \mathbf {H}} (f (x) + g (x)) \leqslant \sup _ {x \in \mathbf {H}} f (x) + \sup _ {x \in \mathbf {H}} g (x).
\]

把第 4 节命题 5 与第 2 节定理 1 应用于此，即得不等式 (17) 与 (18)。其余不等式的证明是类似的。

推论 1. 设 \(f\) 是定义在 \(\mathbf{X}\) 上的实值函数，\(k\) 是实数。则

\[
\sup _ {x \in \mathbf {A}} (f (x) + k) = k + \sup _ {x \in \mathbf {A}} f (x)\tag{22}
\]

只要两边都有定义；又若 \(k \geqslant 0\) ，则

\[
\sup _ {x \in \mathbf {A}} (k f (x)) = k. \sup _ {x \in \mathbf {A}} f (x)\tag{23}
\]

只要两边都有定义。

推论 2. 设 \(f_{1}\) 与 \(f_{2}\) 是分别定义在集合 \(X_{1}, X_{2}\) 上的两个实值函数；则若 \(A_{1}, A_{2}\) 分别是 \(X_{1}, X_{2}\) 的任意非空子集，我们有

\[
\sup _ {(x _ {1}, x _ {2}) \in \mathbf {A} _ {1} \times \mathbf {A} _ {2}} (f _ {1} (x _ {1}) + f _ {2} (x _ {2})) = \sup _ {x _ {1} \in \mathbf {A} _ {1}} f _ {1} (x _ {1}) + \sup _ {x _ {2} \in \mathbf {A} _ {2}} f _ {2} (x _ {2})\tag{24}
\]

只要两边都有定义；又若 \(f_1, f_2\) 都 \(\geqslant 0\) （分别在 \(A_1, A_2\) 中），则我们有

\[
\sup _ {(x _ {1}, x _ {2}) \in \mathbf {A} _ {1} \times \mathbf {A} _ {2}} (f _ {1} (x _ {1}) f _ {2} (x _ {2})) = \sup _ {x _ {1} \in \mathbf {A} _ {1}} f _ {1} (x _ {1}) \sup _ {x _ {2} \in \mathbf {A} _ {2}} f _ {2} (x _ {2}),\tag{25}
\]

只要两边都有定义。

这是前一推论与第 4 节命题 9 的直接推论。

特别地，若 \(\mathbf{A}\) 与 \(\mathbf{B}\) 是 \(\overline{\mathbf{R}}\) 的两个子集，且和的集合 \(\mathbf{A} + \mathbf{B}\) （由和 \(x + y\) ( \(x \in \mathbf{A}, y \in \mathbf{B}\) )组成）有定义，则我们有

\[
\sup (\mathbf {A} + \mathbf {B}) = \sup \mathbf {A} + \sup \mathbf {B}\tag{26}
\]

这里须右端有定义。同样，若 A 与 B 是 \([0, +\infty]\) 的两个子集，则我们有

\[
\sup \mathbf {A B} = \sup \mathbf {A}. \sup \mathbf {B}\tag{27}
\]

只要两边都有定义。

命题 13. 设 \(f\) 与 \(g\) 是定义在滤子化集合 \(X\) 上的两个实值函数。

\begin{enumerate}
\def\labelenumi{(\roman{enumi})}
\tightlist
\item
  我们有
\end{enumerate}

\[
\lim \sup (f + g) \leqslant \lim \sup f + \lim \sup g, \tag {28}
\]

\begin{enumerate}
\def\labelenumi{(\arabic{enumi})}
\setcounter{enumi}{28}
\tightlist
\item
  \(\lim \sup f + \lim \inf g \leqslant \lim \sup (f + g)\)
\end{enumerate}

只要这些不等式两边都有定义。

\begin{enumerate}
\def\labelenumi{(\roman{enumi})}
\setcounter{enumi}{1}
\tightlist
\item
  若 \(f\) 与 \(g\) 都 \(\geqslant 0\) （在 \(\mathbf{X}\) 上），则我们有
\end{enumerate}

\[
\lim \sup f g \leqslant (\lim \sup f) (\lim \sup g), \tag {30}
\]

\[
(\lim \sup f) (\lim \inf g) \leqslant \lim \sup f g \tag {31}
\]

只要这些不等式两边都有定义。

\begin{enumerate}
\def\labelenumi{(\roman{enumi})}
\setcounter{enumi}{2}
\tightlist
\item
  若 \(f \geqslant 0\) （在 \(\mathbf{X}\) 上），则
\end{enumerate}

\[
\lim \sup (1 / f) = 1 / (\lim \inf f)\tag{32}
\]

（约定 \(1 / 0 = +\infty\) ）。

这些关系式是命题 12 与关系式 (10) 的推论。

推论 1. 设 \(f\) 与 \(g\) 是定义在滤子化集合 \(X\) 上的两个实值函数。若 \(\lim g\) 存在，则

\[
\lim \sup (f + g) = \lim \sup f + \lim g \tag {33}
\]

只要两边都有定义，且

\[
\lim \sup f g = (\lim \sup f) (\lim g) \tag {34}
\]

只要两边都有定义且 \(f\) 与 \(g\) 都 \(\geqslant 0\) 。

推论 2. 设 \(f\) 与 \(g\) 是定义在滤子化集合 X 上的两个实值函数。若 \(\lim f = +\infty\) 、\(\lim \inf g > -\infty\) 且 \(f + g\) 有定义，则 \(\lim (f + g) = +\infty\) 。若 \(\lim f = +\infty\) 、\(\lim \inf g > 0\) 且 \(fg\) 有定义，则 \(\lim fg = +\infty\) 。

